\chapter{Preliminary Synthesis Conclusion and Proposed Research Path}

\section{Summary of Hypothetical Findings}
Based on the synthesis of existing NASA flight data and the theoretical framework, this research proposal establishes a foundation for several anticipated conclusions. These findings represent the "state-of-the-pre-simulation" knowledge that defines the research direction.

The following outcomes are anticipated:
\begin{enumerate}
    \item \textbf{Localized Heating Dominance:} The research is expected to confirm that local "scalloping" on the toroidal stack acts as a primary driver of heat flux augmentation. While the global bluntness of the 15m HIAD reduces average heating by $\approx 85\%$, the local crests are predicted to experience spikes that approach or exceed the $14.4 \text{ W/cm}^2$ benchmarks seen in IRVE-3.
    
    \item \textbf{Validation of Recirculation Shielding:} Simulations are expected that toroidal valleys create stable "dead air" vortices. This supports the hypothesis that the valleys will remain within the $400^\circ\text{C}$ structural limit, even as the crests face peak enthalpy.
    
    \item \textbf{Necessity of Real-Gas Modeling:} The research is anticipated to conclude that a \textit{Thermally Perfect Gas} model is mandatory for accuracy. Traditional calorically perfect models, which predict a $T_0$ of 1320 K at Mach 5, likely fail to capture the density variations responsible for the "density hole" observed in IRVE-3 reconstructions.
\end{enumerate}

\section{Addressing the Missing Data Gaps}
The primary objective of the upcoming simulation phase is to generate the high-fidelity data currently missing from the public literature. The research will specifically target:
\begin{itemize}
    \item \textbf{High-Resolution Surface Mapping:} Resolving the exact heat flux and pressure distribution along the "scalloped" wetted length to quantify the heating penalty factor.
    \item \textbf{FTPS Safety Margin Analysis:} Determining the quantitative gap between predicted peak heating and the $103 \text{ W/cm}^2$ limit of Gen-2 SiC materials under scalloped flow conditions.
\end{itemize}

\section{Recommendations for the Simulation Phase}
To ensure the proposed research addresses these gaps, the following technical recommendations are made:

\begin{enumerate}
    \item \textbf{Refine Mesh for SBLI Resolution:} Given the complex Shock-Boundary Layer Interactions (SBLI) predicted in Chapter 4, the mesh must be refined specifically at the reattachment points on the toroidal crests to ensure $y^+ \approx 1$.
    
    \item \textbf{Benchmark Against IRVE-3 Reconstructions:} The initial "Smooth" baseline simulation  should be calibrated to match the IRVE-3 stagnation point data ($14.4 \text{ W/cm}^2$) to ensure the solver is accurately capturing flight-like conditions.
    
    \item \textbf{Sensitivity Study on Gas Constants:} A comparative study between \textit{Ideal Gas} and \textit{Thermally Perfect} assumptions should be conducted to evaluate the impact on shock standoff distance, directly addressing the atmospheric modeling discrepancies found in NASA post-flight analyses.
\end{enumerate}